\documentclass[11pt,a4paper]{article}

\usepackage[margin=1in]{geometry}
\usepackage{amsmath,amssymb,amsthm}
\usepackage{booktabs}
\usepackage{hyperref}

\newcommand{\ebasis}[1]{\mathbf{e}_{#1}}
\newcommand{\conj}[1]{\overline{#1}}
\newcommand{\dOmega}{\,\mathrm{d}\Omega}
\newcommand{\code}[1]{\texttt{#1}}

\title{GSHTrans: mathematics and numerical methods\\
\large A reference for what the library does, and why}
\author{}
\date{}

\begin{document}
\maketitle

\begin{abstract}
This is a record rather than a plan.  Part~I collects the definitions,
conventions and relations the library rests on, in the form it actually
uses them, with a translation table for the three sources whose notations
differ.  Part~II explains the numerical methods: what the algorithms are,
why they were chosen, and the handful of places where the reasoning is
subtle enough to be worth writing down.  Neither part is a derivation.
The companion note \code{canonical-components.tex} carries the
mathematical arguments; \code{core-plan.md} and
\code{field-algebra-plan.md} carry the design decisions and the
measurements behind them.

Statements here are either verified by a test in the library or read
directly from a cited source.  The two places where a convention remains
genuinely open are marked as such.
\end{abstract}

\tableofcontents

\part{Mathematics}

\section{Conventions}
\label{sec:conventions}

\subsection{The canonical basis}

At each point of the sphere, with $\hat{\mathbf{r}},
\hat{\boldsymbol{\theta}}, \hat{\boldsymbol{\phi}}$ the usual orthonormal
spherical polar frame,
\begin{equation}
  \ebasis{\pm} = \mp \tfrac{1}{\sqrt{2}}
    \left( \hat{\boldsymbol{\theta}} \pm i \hat{\boldsymbol{\phi}} \right),
  \qquad
  \ebasis{0} = \hat{\mathbf{r}},
  \qquad \alpha \in \{-1,0,+1\}.
  \label{eq:basis}
\end{equation}
Two consequences are used everywhere:
\begin{equation}
  \conj{\ebasis{\alpha}} = (-1)^{\alpha} \ebasis{-\alpha},
  \qquad
  g_{\alpha\beta} = \ebasis{\alpha}\cdot\ebasis{\beta}
                  = (-1)^{\alpha}\delta_{\alpha+\beta,0}.
  \label{eq:basisconj}
\end{equation}
So $g_{+-}=g_{-+}=-1$, $g_{00}=1$, and $g$ is its own inverse.  The
library works with contravariant components throughout and applies $g$
explicitly wherever a contraction is taken; index position is not tracked.

Under a rotation of the tangent frame through $\psi$ about
$\hat{\mathbf{r}}$, $\ebasis{\pm}\mapsto e^{\mp i\psi}\ebasis{\pm}$.

\subsection{What the library stores, exactly}

The single most useful statement for reading the code:

\begin{quote}
  \code{Wigner}'s stored value at upper index $N$, degree $l$ and order
  $m$ is
  \begin{equation}
    \left(\frac{2l+1}{4\pi}\right)^{1/2} d^{l}_{Nm}(\theta)
      \;=\; X^{N}_{lm}(\theta),
    \label{eq:stored}
  \end{equation}
  which is exactly the scalar harmonic $X^N_{lm}$ of Dahlen \& Tromp
  (1998), their (C.117), built on their generalized Legendre function
  $P^N_{lm} = d^l_{Nm}$ --- \emph{upper index first}.
\end{quote}

This is verified rather than assumed.  \code{tests/CheckWignerConvention.h}
compares all nine $l=1$ values against D\&T (C.115); the check
discriminates, because the four entries with $N-m$ odd distinguish
$d^l_{Nm}$ from $d^l_{mN}$, which differ by $(-1)^{N-m}$.  Two further
facts corroborate it: the recursion seed evaluates to D\&T (C.125), and
\code{tests/CheckLegendre.h} pins the $N=0$ row against
\code{std::sph\_legendre}, fixing the orthonormal scaling and the
Condon--Shortley phase.

\subsection{Translation table}

The three standard sources use different notations, and two of them differ
from this library in normalisation.  Everything in this table is read
directly from the cited equations.

\begin{center}
\begin{tabular}{llll}
  \toprule
  & harmonic & colatitudinal part & normalised? \\
  \midrule
  GSHTrans
    & $Y^{N}_{lm} = X^{N}_{lm} e^{im\phi}$
    & $X^{N}_{lm} = \left(\tfrac{2l+1}{4\pi}\right)^{1/2} d^{l}_{Nm}$
    & yes \\
  Dahlen \& Tromp (C.117)
    & $Y^{N}_{lm} = X^{N}_{lm} e^{im\phi}$
    & $X^{N}_{lm} = \left(\tfrac{2l+1}{4\pi}\right)^{1/2} P^{N}_{lm}$
    & yes \\
  Woodhouse \& Deuss [41]
    & $Y^{Nm}_{l} = P^{Nm}_{l} e^{im\phi}$
    & $P^{Nm}_{l} = d^{(l)}_{Nm}$
    & \textbf{no} \\
  Phinney \& Burridge (1.5)
    & $X^{Nm} = P^{Nm}_{l} e^{im\phi}$
    & $P^{Nm}_{l}$, their Appendix A
    & \textbf{no} \\
  \bottomrule
\end{tabular}
\end{center}

Three things to take from it.

\paragraph{The index order is the same in all four.}  D\&T write
$P^{N}_{lm}$, Woodhouse and P\&B write $P^{Nm}_{l}$, and this library
writes $d^{l}_{Nm}$; all mean the same function with the upper index
first.  D\&T state that they follow P\&B in the definition.

\paragraph{The normalisation is not.}  D\&T and this library carry the
$[(2l+1)/4\pi]^{1/2}$; Woodhouse and P\&B do not.  A formula taken from
Woodhouse or P\&B therefore needs that factor supplied, and a formula
taken from D\&T does not.

\paragraph{The transposed function is a different thing.}  D\&T (C.118)
gives $P^{m}_{lN} = (-1)^{m+N} P^{N}_{lm}$, so getting the index order
wrong costs a sign whenever $m-N$ is odd, and nothing when it is even ---
which is exactly the kind of error that survives half a test suite.

\subsection{What remains a convention}
\label{sec:remains}

This document used to list two statements as unsettled: the sign $\mp$
in \eqref{eq:basis}, since some authors use
$\ebasis{\pm}=\pm(\hat{\boldsymbol{\theta}}\pm
i\hat{\boldsymbol{\phi}})/\sqrt{2}$, which flips every odd-rank reality
phase; and the overall signs of $\eth$ and $\conj{\eth}$
(\S\ref{sec:eth}).  Each was said to be unobservable, and each was, for
the same reason: everything the library computed stayed in canonical
components from end to end, and a convention that is never read back
into the physical frame cannot be seen.

They are now pinned, jointly, by reading a gradient back into the
physical frame and requiring it to be the gradient.  Inverting
\eqref{eq:basis} for a vector $\mathbf{v}=v^{\alpha}\ebasis{\alpha}$,
\begin{equation}
  v_{\theta} = \frac{v^{-}-v^{+}}{\sqrt{2}},
  \qquad
  v_{\phi} = -i\,\frac{v^{+}+v^{-}}{\sqrt{2}},
  \qquad
  v_{r} = v^{0},
  \label{eq:physical}
\end{equation}
and \code{tests/TestConventions.cpp} requires that for a scalar $f$ the
surface gradient satisfies
$\mathbf{v}=\hat{\boldsymbol{\theta}}\,\partial_{\theta}f
+\hat{\boldsymbol{\phi}}\,(\sin\theta)^{-1}\partial_{\phi}f$ pointwise,
and that the metric trace of the surface gradient of a tangential vector
field is its surface divergence.  Both hold to $10^{-13}$, and both fail
by an $O(1)$ amount if either sign is flipped alone.

What is left is genuinely a convention and nothing more: flipping
\eqref{eq:basis} \emph{and} the sign of the $\eth$ pair together is a
relabelling of which component is called $+1$, and \eqref{eq:physical}
flips with it.  So there is one convention here rather than two unknowns,
this document states it, and a test observes it.

\paragraph{A trap worth recording.}  A tangential field written as
$a(\theta,\phi)\,\hat{\boldsymbol{\theta}}$ is almost never smooth, because
$\hat{\boldsymbol{\theta}}$ is not: at a pole its direction depends on the
azimuth from which the pole is approached.  The innocuous-looking
$a=\sin\theta\cos\phi$ is not differentiable there, its spin-weighted
expansion does not converge, and the divergence check above came out wrong
by $0.18$ at every truncation from $l_{\max}=8$ to $128$.  That is what a
non-convergent expansion looks like, and it is easily mistaken for a bug in
the operator.  The test fields are built as
$\nabla_1 f + \hat{\mathbf{r}}\times\nabla_1 g$ from smooth scalars, which
is smooth by construction.

\section{Canonical components}
\label{sec:components}

A rank-$p$ tensor field has $3^p$ scalar components, each labelled by a
multi-index $(\alpha_1\ldots\alpha_p)\in\{-1,0,1\}^p$.  Under the frame
rotation, the component $T^{\alpha_1\ldots\alpha_p}$ acquires the phase
$e^{-iN\psi}$ with
\begin{equation}
  N = \sum_{i=1}^{p} \alpha_i,
  \label{eq:N}
\end{equation}
and it is this $N$ --- not the multi-index --- that is the upper index of
the component's harmonic expansion.

\paragraph{The multi-index is not the upper index.}  For a vector they
coincide, $N=\alpha$, which is the coincidence that misleads.  For
$p\geq 2$ distinct components share an upper index: a rank-2 tensor has
three at $N=0$, namely $(-+)$, $(00)$ and $(+-)$.  A collection of
components labelled only by $N$ therefore does not determine a tensor.
The number carrying a given $N$ is the trinomial coefficient; for $p=2$
these are $1,2,3,2,1$ across $N=-2\ldots 2$.

\paragraph{Products and contractions.}  Concatenating multi-indices adds
upper indices by \eqref{eq:N}, which is the statement ``upper indices add
under pointwise multiplication''.  Contracting slots $j$ and $k$ against
the metric gives
\begin{equation}
  \left(\mathrm{tr}_{jk} T\right)^{\ldots}
    = \sum_{\alpha}(-1)^{\alpha}\,T^{\ldots\,\alpha\,\ldots\,-\alpha\,\ldots}
    = -T^{\ldots -\ldots +\ldots} + T^{\ldots 0\ldots 0\ldots}
      - T^{\ldots +\ldots -\ldots},
  \label{eq:contract}
\end{equation}
the contracted pair contributing $\alpha+(-\alpha)=0$, so the result
carries the upper index of the surviving slots.

\section{Generalized spherical harmonics}
\label{sec:gsh}

A field $f$ of upper index $N$ expands as
\begin{equation}
  f(\theta,\phi) = \sum_{l\geq|N|}\sum_{m=-l}^{l} f^{N}_{lm} Y^{N}_{lm},
  \qquad
  f^{N}_{lm} = \int_{S^2}\conj{Y^{N}_{lm}}\, f \dOmega,
  \label{eq:expansion}
\end{equation}
with
\begin{equation}
  Y^{N}_{lm}(\theta,\phi)
    = \left(\frac{2l+1}{4\pi}\right)^{1/2} d^{l}_{Nm}(\theta)\, e^{im\phi}.
  \label{eq:gsh}
\end{equation}
These are orthonormal, $\int \conj{Y^N_{lm}} Y^N_{l'm'} \dOmega =
\delta_{ll'}\delta_{mm'}$.  The sum starts at $l=|N|$ because $d^l_{Nm}$
vanishes identically below it: there is no harmonic of that upper index
at lower degree, and correspondingly no coefficient to store.  This is
the single fact most likely to look like a bug on first contact --- see
\S\ref{sec:projection}.

\paragraph{The normalisation is pinned by the code.}  Taking $f$ constant,
\eqref{eq:expansion} gives $f^0_{00} = 2\sqrt{\pi}f$, and
\code{GaussLegendreGrid} implements exactly that on its one-point grid.

\paragraph{Conjugation.}  Using $d^l_{-m,-N}=(-1)^{m-N}d^l_{mN}$ and the
reality of $d$, equation \eqref{eq:gsh} gives D\&T (C.109),
\begin{equation}
  \conj{Y^{N}_{lm}} = (-1)^{m+N} Y^{-N}_{l,-m},
  \label{eq:Yconj}
\end{equation}
from which follow the two spectral reality relations of
\S\ref{sec:reality}.

\section{The generalized Legendre functions}
\label{sec:legendre}

$P^N_{lm}(\mu)=d^l_{Nm}(\theta)$, $\mu=\cos\theta$, is the regular
solution of the generalized Legendre equation, normalised as D\&T (C.112)
so that the addition theorem below holds.  It does \emph{not} reduce to
the associated Legendre function at $N=0$; D\&T (C.113) relates them.

\subsection{Symmetries}

D\&T (C.118), in full, since the library uses two of them and could use
two more:
\begin{align}
  P^{N}_{l,-m}(\mu) &= (-1)^{l+N} P^{N}_{lm}(-\mu), &
  P^{-N}_{lm}(\mu)  &= (-1)^{l+m} P^{N}_{lm}(-\mu), \notag \\
  P^{-N}_{l,-m}(\mu) &= (-1)^{m+N} P^{N}_{lm}(\mu), &
  P^{m}_{lN}(\mu)   &= (-1)^{m+N} P^{N}_{lm}(\mu),
  \label{eq:symmetries}\\
  P^{-m}_{l,-N}(\mu) &= P^{N}_{lm}(\mu). \notag
\end{align}
The third is the $(m,N)\mapsto(-m,-N)$ involution; the second is the
$\theta\mapsto\pi-\theta$ reflection, which maps the Gauss--Legendre node
set to itself exactly, since those nodes are symmetric about $\pi/2$.
Composed, the two generate a group of order four acting on
$(N,m,\theta)$, so the table of stored values could in principle be
reduced fourfold.  Neither reduction is implemented; \S\ref{sec:whatnext}
says why.  The fourth relation is the transposition that the convention
check discriminates against.

\subsection{Boundary values}

The values at the extreme orders have closed forms, D\&T (C.123) and
(C.125), and these are precisely the library's
\code{WignerMinOrder} and \code{WignerMaxOrder}:
\begin{align}
  P^{N}_{l,-l} &= \left[\frac{(2l)!}{(l+N)!(l-N)!}\right]^{1/2}
                  \left(\sin\tfrac{\theta}{2}\right)^{l+N}
                  \left(\cos\tfrac{\theta}{2}\right)^{l-N},
  \label{eq:minorder}\\
  P^{N}_{ll}   &= (-1)^{l+N} P^{N}_{l,-l}\big|_{N\to-N}.
  \label{eq:maxorder}
\end{align}
They are the seeds of D\&T's recursion in $m$, and the boundary terms of
the library's recursion in $l$ (\S\ref{sec:recursion}).

\subsection{The addition theorem}

D\&T (C.127),
\begin{equation}
  \sum_{m=-l}^{l} P^{N}_{lm}(\mu)\, P^{N'}_{lm}(\mu) = \delta_{NN'}.
  \label{eq:addition}
\end{equation}
This is the strongest available check on a computed table, and
\code{tests/CheckAdditionTheorem.h} is it.  Because the library stores
$X^N_{lm}$ rather than $P^N_{lm}$, the right-hand side there is
$\delta_{NN'}(2l+1)/4\pi$.

\section{Reality}
\label{sec:reality}

For a real tensor field, applying \eqref{eq:basisconj} to each factor
gives
\begin{equation}
  T^{-\alpha_1\ldots-\alpha_p} = (-1)^{N}\,
    \conj{T^{\alpha_1\ldots\alpha_p}}.
  \label{eq:reality}
\end{equation}
For a vector this reads $u^0=\conj{u^0}$ and $u^-=-\conj{u^+}$: the radial
component is a real field and the two transverse components are not
independent.

\paragraph{Which components must be stored.}  \eqref{eq:reality} pairs
each multi-index with its negation, and one representative of each orbit
suffices.  The only fixed point is the all-zero index, which is therefore
a real-valued field; every other orbit has size two.  With a permutation
symmetry present the orbits are those of the group generated by the
permutations \emph{together with} the negation, and a component whose
orbit contains its own negation is pinned to a single real number ---
real, or purely imaginary when a permutation sign intervenes.  Such
self-paired components always sit at $N=0$: permutation preserves the slot
sum and negation reverses it, so a component fixed by the combination
satisfies $N=-N$.

The resulting counts, computed by the library rather than tabulated by
hand, are the independent real degrees of freedom in every case:

\begin{center}
\begin{tabular}{lrrr}
  \toprule
  & complex components & pinned & reals per point \\
  \midrule
  rank 1                & 1  & 1 & 3 \\
  rank 2                & 4  & 1 & 9 \\
  rank 2, symmetric     & 2  & 2 & 6 \\
  rank 2, antisymmetric & 1  & 1 & 3 \\
  rank 3, symmetric     & 4  & 2 & 10 \\
  rank 4                & 40 & 1 & 81 \\
  rank 4, elastic       & 8  & 5 & 21 \\
  \bottomrule
\end{tabular}
\end{center}

The last row is worth noting: nothing in the machinery knows about
elasticity, and $c_{ijkl}=c_{jikl}=c_{ijlk}=c_{klij}$ comes out at the
familiar 21.

\paragraph{In the spectral domain} the same relation takes two distinct
forms, and it is worth keeping them apart.  For \emph{any} field $f$ of
upper index $N$, \eqref{eq:Yconj} gives the coefficients of $\conj{f}$,
which has upper index $-N$:
\begin{equation}
  \left(\conj{f}\right)^{-N}_{l,-m} = (-1)^{m-N}\conj{f^{N}_{lm}}.
  \label{eq:basiclevel}
\end{equation}
For a negation-paired component of a \emph{real tensor}, \eqref{eq:reality}
supplies a further $(-1)^N$ and the two combine to
\begin{equation}
  T^{-N}_{l,-m} = (-1)^{m}\conj{T^{N}_{lm}},
  \label{eq:complevel}
\end{equation}
which at $N=0$ is the familiar condition on a real scalar field, and is
what the reduced $m\geq 0$ coefficient storage implements.

\section{Raising and lowering}
\label{sec:eth}

$\eth$ and $\conj{\eth}$ change the upper index by one.  On a basis
function,
\begin{equation}
  \eth\, Y^{N}_{lm} = -\sqrt{(l-N)(l+N+1)}\; Y^{N+1}_{lm},
  \qquad
  \conj{\eth}\, Y^{N}_{lm} = +\sqrt{(l+N)(l-N+1)}\; Y^{N-1}_{lm},
  \label{eq:eth}
\end{equation}
so in the spectral domain each is a multiplication by an $l$-dependent
factor, with no coupling between different $(l,m)$.  They are the
operations that connect different upper indices, and they are diagonal in
$(l,m)$ rather than local in $(\theta,\phi)$ --- which is why the library
has two representations rather than one, and why an expression such as
``the gradient of a product'' is necessarily evaluated in both.

The surface gradient of a scalar has canonical components proportional to
$\eth f$ and $\conj{\eth} f$, at upper index $+1$ and $-1$, with vanishing
radial component.

\paragraph{Two identities, and a third that pins the rest.}  Composing
the factors,
\begin{equation}
  \conj{\eth}\eth\big|_{N=0} = -l(l+1),
  \qquad
  [\eth,\conj{\eth}] = -2N.
  \label{eq:ethidentities}
\end{equation}
The first says $\conj{\eth}\eth$ is the surface Laplacian on a scalar; it
pins the relative sign of the two operators and both magnitudes.  The
second is structural.  Both are tested, and both survive flipping the sign
of the pair --- which is why this document long carried the overall sign
as an open question for Phinney \& Burridge.

It does not need one.  P\&B do not use $\eth$ as such, and the sign is not
theirs to settle; what has to be right is that gradients of functions come
out as gradients, which is checkable here.  Requiring exactly that, in the
physical frame, fixes the overall sign in \eqref{eq:eth} against the sign
in \eqref{eq:basis}: see \S\ref{sec:remains}.  The two are one convention,
stated in this document and observed by
\code{tests/TestConventions.cpp}.

\paragraph{The degree ranges look after themselves.}  Raising a field at
$N\geq 0$ produces one starting at $l=N+1$, and the factor at $l=N$ is
$\sqrt{0}=0$: the coefficient with nowhere to go was going to vanish.
Lowering produces a field starting at a lower degree, whose extra
coefficient has no source and is zero, which is what having no content
there means.

\section{The contravariant derivative, and how it differs}
\label{sec:contravariant}

Phinney \& Burridge, and Dahlen \& Tromp following them, work not with
$\eth$ but with a \emph{contravariant derivative} $\partial^\sigma$,
$\sigma\in\{-,0,+\}$, which takes the expansion coefficients of a rank-$q$
tensor to those of the rank-$(q+1)$ tensor $\nabla\mathbf{T}$.  It is
closely related to $\eth$ and it is not the same operator, and the
difference matters for anyone using this library on tensors.

\subsection{The definition}

With D\&T (C.146),
\begin{equation}
  \Omega^{\pm N}_{l} = \sqrt{\tfrac{1}{2}(l \pm N)(l \mp N + 1)},
  \label{eq:omega}
\end{equation}
their (C.151)--(C.153) read
\begin{align}
  \partial^{-} T^{\alpha_1\ldots\alpha_q}_{lm}
    &= r^{-1}\Big[\Omega^{+N}_{l}\, T^{\alpha_1\alpha_2\ldots\alpha_q}_{lm}
      - T^{(\alpha_1-1)\alpha_2\ldots\alpha_q}_{lm}
      - \cdots
      - T^{\alpha_1\alpha_2\ldots(\alpha_q-1)}_{lm}\Big],
    \label{eq:dminus}\\
  \partial^{0} T^{\alpha_1\ldots\alpha_q}_{lm}
    &= \frac{d\,T^{\alpha_1\ldots\alpha_q}_{lm}}{dr},
    \label{eq:dzero}\\
  \partial^{+} T^{\alpha_1\ldots\alpha_q}_{lm}
    &= r^{-1}\Big[\Omega^{-N}_{l}\, T^{\alpha_1\alpha_2\ldots\alpha_q}_{lm}
      - T^{(\alpha_1+1)\alpha_2\ldots\alpha_q}_{lm}
      - \cdots
      - T^{\alpha_1\alpha_2\ldots(\alpha_q+1)}_{lm}\Big],
    \label{eq:dplus}
\end{align}
with the stipulation that any coefficient whose shifted index leaves
$\{-1,0,1\}$ is zero.  The components of the gradient are
$(\nabla\mathbf{T})^{\sigma\alpha_1\ldots\alpha_q} = \partial^{\sigma}
T^{\alpha_1\ldots\alpha_q}$: the operator \emph{prepends a slot}, and the
result carries upper index $\sigma + N$, which is \eqref{eq:N} applied to
the lengthened multi-index.

\subsection{Two differences, one benign and one not}

\paragraph{A factor of $\sqrt2$, which is only the basis normalisation.}
Comparing \eqref{eq:omega} with \eqref{eq:eth},
\begin{equation}
  \left|\text{raising factor}\right| = \sqrt{(l-N)(l+N+1)}
    = \sqrt{2}\,\Omega^{-N}_{l},
  \qquad
  \left|\text{lowering factor}\right| = \sqrt{2}\,\Omega^{+N}_{l}.
  \label{eq:ethomega}
\end{equation}
The $\sqrt2$ is exactly the $1/\sqrt2$ in $\ebasis{\pm} = \mp(
\hat{\boldsymbol{\theta}} \pm i\hat{\boldsymbol{\phi}})/\sqrt2$, which
appears explicitly in D\&T's surface gradient (C.144).  So on the unit
sphere, up to that factor and to the sign convention of
\S\ref{sec:eth}, $\partial^{+}$ is $\eth$ and $\partial^{-}$ is
$\conj{\eth}$ --- \emph{for a scalar}.

\paragraph{Connection terms, which are not benign.}  For $q\geq 1$ the
bracket in \eqref{eq:dminus} and \eqref{eq:dplus} is not a multiplication.
It subtracts the components with one slot index shifted, once per slot.
Those terms are the connection: the canonical basis vectors themselves
vary over the sphere,
\begin{equation}
  \left[\pm\partial_\theta + i(\sin\theta)^{-1}\partial_\phi\right]
    \ebasis{\alpha}
    = \alpha\cot\theta\,\ebasis{\alpha} - \sqrt{2}\,\ebasis{\alpha\pm1},
  \label{eq:basisderiv}
\end{equation}
so differentiating a tensor field differentiates its basis too.  A scalar
has no slots and therefore no connection terms, which is why the two
operators coincide there and only there.

\subsection{The consequence for this library}

\begin{quote}
  \textbf{Applying $\eth$ to each component of a tensor does not give the
  gradient of the tensor.}  Use \code{SurfaceGradient}.  It gives the $\eth$ of each component as a
  spin-weighted function, which for $q\geq 1$ is not a component of
  $\nabla\mathbf{T}$ and in general is not a component of anything.
\end{quote}

This is a trap rather than a limitation: \code{Raise} and \code{Lower} are
correct for what they claim --- a single spin-weighted function --- and are
correct for the surface gradient of a scalar, which is the case most people
reach for first.  They are simply not the \emph{ambient} tensor operator, and
nothing in their types says so, because a tensor expansion's components are
ordinary spin expansions.

\subsection{The intrinsic derivative, where \texorpdfstring{$\eth$}{eth} is
exactly right}
\label{sec:intrinsic}

The paragraph above must not be read as forbidding component-wise $\eth$
altogether, because there is a case in which it is precisely the correct
operator.

There are two derivatives on the sphere.  $\partial^\sigma$ is the
\emph{ambient} one: it differentiates the tensor and its basis, and the basis
leaves the tangent plane.  The \emph{intrinsic} covariant derivative --- the
Levi-Civita connection of the induced metric, which is what surface
differential geometry means by differentiation --- is closed on
\emph{tangential} tensors, those whose every slot is $\pm1$ with no radial
index.

On such a tensor the connection terms vanish identically.  Every shift
$\alpha_i + \sigma$ with $\alpha_i,\sigma \in \{\pm1\}$ lands on $-2$, $0$ or
$+2$; the outer two leave the alphabet, and $0$ names a component with a
radial slot, which a tangential tensor does not have.  What is left is pure
$\Omega$ multiplication.  Splitting $\nabla_1 T$ by whether the inherited slot
is radial,
\begin{align}
  (\nabla_1 T)^{\sigma \alpha_1\ldots\alpha_q}
    &= \Omega^{\mp N}_{l}\, T^{\alpha_1\ldots\alpha_q},
    &&\text{all } \alpha_i \text{ tangential},
    \label{eq:intrinsicpart}\\
  (\nabla_1 T)^{\sigma \alpha_1\ldots 0_j \ldots\alpha_q}
    &= -\,T^{\alpha_1\ldots \sigma_j \ldots\alpha_q},
    &&\text{one slot radial}.
    \label{eq:curvature}
\end{align}
Both hold \emph{exactly}, and are checked as such against
\code{SurfaceGradient} applied to a rank-1 tensor with its radial component
set to zero.

\eqref{eq:curvature} is the extrinsic curvature: on a unit sphere the normal
part of $\nabla_a v_b$ is exactly $-v$.  \eqref{eq:intrinsicpart} says

\begin{quote}
  \textbf{On a tangential tensor the intrinsic covariant derivative is $\eth$
  applied component by component, up to $\sqrt2$, and it is closed.}
\end{quote}

So the two operators differ by an algebraic term, not a differential one, and
$\eth$ is not a poor relation of $\partial^\sigma$: it is the intrinsic
derivative, exact on exactly the objects --- tangential, spin-weighted --- that
it was invented for.  Which of the two a caller wants depends on whether the
radial directions are part of the problem.

\subsection{What is implemented}

\code{SurfaceGradient} is $\nabla_1$: \eqref{eq:dminus} and
\eqref{eq:dplus} without the $r^{-1}$, taking a rank-$q$ tensor expansion
to a rank-$(q+1)$ one.

\paragraph{No radial part.}  $\partial^0$ is $\partial/\partial r$, which
an angular library cannot supply; but the \emph{surface} gradient has no
$\ebasis{0}$ component at all (D\&T C.145), so the operator is complete
without it and the radial derivative is the user's to provide.  What that
means for the result is that its $\sigma=0$ block is stored and zero: a
rank-$(q+1)$ tensor whose $\ebasis{0}$ slot vanishes is an ordinary
tensor, and keeping it as one leaves the result composable with
contraction, further gradients and the transform, for a third of the
storage.

\paragraph{The result has no symmetry and keeps the operand's reality.}
$\nabla_1$ of a symmetric tensor is not symmetric in the new slot against
the old ones, and inferring a symmetry from an operator is the problem
\code{Materialise} declined to solve.  $\nabla_1$ of a real tensor is
real, so only the reduced set is computed.

\paragraph{It needed one thing phase 5 had left out.}  Every term is a
coefficient of some component at fixed $(l,m)$, and the shifted components
need not be stored.  \code{TensorExpansion::Coefficient} now answers for
\emph{any} representable component, resolving the stored case, a
permutation relative, a reality relative by \eqref{eq:complevel}, a real
block's reduced $m\geq 0$ storage, a pinning as imaginary, and a vanishing
orbit.  It is the spectral counterpart of phase~4's component accessor and
closes the asymmetry noted above.

\paragraph{What the tests pin.}  On a scalar the operator is $\eth$ up to
$\sqrt2$, which checks the $\Omega$ factors and nothing else.  The metric
trace of the \emph{second} gradient of a scalar is the surface Laplacian,
$-l(l+1)$ --- and that runs through the connection terms, since
$\partial^-$ of the $+1$ component subtracts the $0$ component, so the
identity fails outright if the operator is $\eth$ applied component-wise.
The chain rule \eqref{eq:leibniz} is checked at rank one, in the spatial
domain.  And the gradient of a real tensor is checked against the same
field widened to a complex one, which exercises reading derived components
and writing pinned ones.

\paragraph{Divergence and curl} are contractions of the gradient (D\&T
C.6.2).  Contraction is pointwise, so a caller evaluates the gradient and
contracts with the machinery of \S\ref{sec:types}; there is no separate
operator.

\paragraph{One further property worth knowing.}  $\partial^{\sigma}$
preserves the chain rule as an operator on fields, D\&T (C.154)--(C.155):
$\nabla(\mathbf{T}\mathbf{P}) = (\nabla\mathbf{T})\mathbf{P} +
\mathbf{T}(\nabla\mathbf{P})$ gives
\begin{equation}
  \partial^\sigma(T^{\alpha\ldots}P^{\beta\ldots}) =
  (\partial^\sigma T^{\alpha\ldots})P^{\beta\ldots} +
  T^{\alpha\ldots}(\partial^\sigma P^{\beta\ldots}).
  \label{eq:leibniz}
\end{equation}  It does \emph{not}
hold coefficient by coefficient, because a product of fields is not a
product of coefficients --- which is the same fact that makes ``the
gradient of a product'' an operation in two representations.

\section{Quadrature}
\label{sec:quadrature}

\paragraph{Colatitude.}  Gauss--Legendre with $n_\theta = l_{\max}+1$
nodes, which integrates a polynomial of degree $2n_\theta-1$ in
$\cos\theta$ exactly.  The nodes are the roots of $P_{n_\theta}$, mapped
by $\theta = \arccos(-x)$, and are symmetric about $\pi/2$.

\paragraph{Longitude.}  The trapezoid rule on $n_\phi$ equally spaced
points, which is exact for $e^{i(m-m')\phi}$ only when $|m-m'| < n_\phi$.
Analysing a band-$l_{\max}$ field needs $|m|,|m'|\leq l_{\max}$, hence
\begin{equation}
  n_\phi \geq 2 l_{\max} + 1.
  \label{eq:nphi}
\end{equation}
At $n_\phi = 2l_{\max}$ the orders $m=+l_{\max}$ and $m=-l_{\max}$ are the
same discrete mode and cannot be separated.  The library takes the
smallest \emph{fast} FFT length at or above the bound --- restricted to
$2^a 3^b 5^c 7^d 11^e 13^f$ with $e+f\leq 1$, the set FFTW has codelets
for --- because the naive $2l_{\max}+2$ lands on $514 = 2\times 257$ at
$l_{\max}=256$, and 257 is prime.  The smooth choice there is $520$, at
about 1.5\% more FFT work, and the FFT is not the expensive stage
(\S\ref{sec:transform}).

\paragraph{Headroom.}  A product of two band-$L$ fields has band $2L$, and
$\int|f|^2$ for band-$L$ $f$ integrates a degree-$2L$ quantity.  Both want
a grid resolving more than the fields do.  \code{Grid::ForBand(L, nMax,
oversampling)} expresses that distinction --- the band of the fields
against the resolution of the grid --- rather than leaving a caller to
compute a degree and hope; the $3/2$ dealiasing rule is
$\text{oversampling}=1.5$ and $2$ is exact for a single product.

\part{Numerical methods}

\section{The transform}
\label{sec:transform}

\subsection{The algorithm}

The forward transform of a field at upper index $n$ is
\begin{equation}
  f^{n}_{lm} = \sum_{i} w_i\, X^{n}_{lm}(\theta_i)\, F_m(\theta_i),
  \qquad
  F_m(\theta) = \frac{2\pi}{n_\phi}\sum_{j} f(\theta,\phi_j) e^{-im\phi_j},
  \label{eq:forward}
\end{equation}
so it is an FFT in longitude at each colatitude, followed by a weighted
sum over colatitudes against the stored $X$.  The inverse runs the same
two stages in the other order.  Both are separable exactly because
\eqref{eq:gsh} is a product of a function of $\theta$ and $e^{im\phi}$.

\paragraph{Cost.}  The FFT stage is $O(n_\theta\, n_\phi \log n_\phi)$ and
the Legendre stage is $O(n_\theta\, l_{\max}^2)$, and at production sizes
the second dominates completely.  Measured at $l_{\max}=256$, $n=2$: a
forward transform takes 19.0\,ms, of which the FFT stage is
\textbf{0.43\,ms} and the Legendre stage \textbf{18.6\,ms} --- a factor of
44.  Any optimisation effort that does not address the Legendre stage is
misdirected, and this number is the reason the library's performance work
concentrates entirely there.

\paragraph{Where the Legendre stage's time actually goes} is less obvious
than it looks.  Single-threaded it runs at about 19\% of the read
bandwidth one core can achieve, so it is not DRAM-bound; what binds it is
load/store throughput.  Per stored value the inner loop does one 8-byte
read that reaches memory and then, from cache, a read of the FFT bin and a
read-modify-write of the coefficient --- four memory operations, only one
of which is a miss.  Three hand-optimised rewrites of the kernel (weight
hoisted, real and imaginary parts split, raw pointers) moved it by less
than a factor of two in either direction.

\subsection{Real fields}

A real-valued field exists only at $n=0$, where $f_{l,-m} = (-1)^m
\conj{f_{lm}}$ by \eqref{eq:complevel}.  The transform then uses a
real-to-complex FFT and stores only $m\geq 0$, halving both the FFT work
and the coefficient count.  At $n\neq 0$ the same reduction is
\emph{invalid}: it assumes the self-relation $f^N_{l,-m}=(-1)^{m-N}
\conj{f^N_{lm}}$, which holds only when $f$ is its own conjugate, and a
spin-weighted field is its own conjugate only at $N=0$.  The library
therefore forbids real-valued fields at nonzero upper index in the type
system, and the transform throws if asked.

\subsection{Transforming is a projection}
\label{sec:projection}

A component at upper index $N$ has no content below degree $|N|$, because
no harmonic there exists.  So the forward transform of an \emph{arbitrary}
field discards whatever that field had at low degree, and evaluating the
result does not reproduce the input.  It reproduces its projection onto
the span of the harmonics that exist.

This is not an error and not a loss of accuracy: it is what the expansion
means.  Once a field is band-limited --- for instance because it came from
an inverse transform --- the round trip is an identity to rounding.
Example~06 shows both, and prints $5\times10^{-1}$ for the first and
$4\times10^{-14}$ for the second.

\section{Computing the \texorpdfstring{$d$}{d}-functions}
\label{sec:recursion}

This is where nearly all the arithmetic is, and where the library differs
most from the reference texts.

\subsection{Recursion in \texorpdfstring{$l$}{l}, not in \texorpdfstring{$m$}{m}}

D\&T describe two schemes, both at fixed degree: their (C.121) recurs in
$m$ at fixed $N$, and (C.122) recurs in $N$ at fixed $m$.  For stability
they iterate (C.121) upward from $m=-l$ and downward from $m=+l$ and meet
at $m = N\cos\theta$, seeding each end with the closed forms
\eqref{eq:minorder}--\eqref{eq:maxorder}.

The library recurs in \emph{degree}, at fixed $(N,m)$.  This is the scheme
Woodhouse used, and the code is an evolution of his Fortran~77
implementation; D\&T do not discuss it.  The recursion is
\begin{equation}
  d^{l}_{Nm} = f_1\, d^{l-1}_{Nm} - f_2\, d^{l-2}_{Nm},
  \label{eq:lrecursion}
\end{equation}
with
\begin{equation}
  f_1 = \frac{(2l-1)\left[\,l(l-1)\cos\theta - Nm\,\right]}
             {(l-1)\sqrt{(l^2-N^2)(l^2-m^2)}},
  \qquad
  f_2 = \frac{l\,\sqrt{\left((l-1)^2-N^2\right)\left((l-1)^2-m^2\right)}}
             {(l-1)\sqrt{(l^2-N^2)(l^2-m^2)}} .
  \label{eq:lcoefficients}
\end{equation}

\paragraph{Why this one.}  The transform consumes the values one degree at
a time, in ascending order, for a fixed $(N,\theta)$ --- see
\eqref{eq:forward}.  A recursion in $l$ produces them in exactly that
order, so the same loop that fills the table can in principle feed the
transform directly, and a truncated transform can stop early and pay for
nothing it does not use.  A recursion in $m$ produces a whole degree at
once and would have to be run to completion first.

\paragraph{What it costs.}  Of the coefficients in
\eqref{eq:lcoefficients}, only the $\cos\theta$ term depends on the
colatitude; the four square roots are $\theta$-independent and are read
from a precomputed table of $\sqrt{k}$ and $1/\sqrt{k}$ --- two vectors of
$2l_{\max}+1$ entries, which is the whole of the auxiliary storage the
recursion needs.  The inner loop runs over $m$ at fixed $(l,N,\theta)$
with unit stride on three rows, which vectorises.

\paragraph{The removable singularity.}  At $l = |N|+1$ the factor $1/(l-1)$
in \eqref{eq:lcoefficients} is $1/|N|$, which is a division by zero at
$N=0$; but $N/(l-1) = \mathrm{sign}(N)$ there, and $d^{l-2}$ does not
exist, so that degree is taken by a separate one-term branch with the
singular factor cancelled analytically.  The same one-term form appears at
the ``critical degree'' $l = m_{\max}+1$, where the row stops growing.

\subsection{The seed row and the boundary terms}

The recursion needs starting values in two senses.

\paragraph{The seed row}, at $l=|N|$, is a loop over $m$ at a single
degree.  At $l=|N|$ the closed form reduces to
\begin{equation}
  P^{N}_{l,-m}\big|_{l=|N|}
    = \sqrt{\binom{2l}{l+m}}\;
      \left(\sin\tfrac{\theta}{2}\right)^{l-m}
      \left(\cos\tfrac{\theta}{2}\right)^{l+m},
  \label{eq:seedrow}
\end{equation}
and the binomial row is generated exactly by the standard step
$\binom{2l}{k} = \binom{2l}{k-1}(2l-k+1)/k$.  The row is at most $2|N|+1$
long --- five entries for a rank-2 application --- so this is not about
cost; see \S\ref{sec:signgam}.

\paragraph{The boundary terms} are the values at $m=\pm l$, needed once
per degree as the row grows.  Evaluating \eqref{eq:minorder} directly costs
three $\log\Gamma$ and an exponential each time, which at $l_{\max}=256$
is about $1.3\times10^5$ transcendental evaluations per upper index per
pass.  They obey a one-term recursion in $l$: writing $s=\sin(\theta/2)$,
$c=\cos(\theta/2)$,
\begin{equation}
  \frac{P^{N}_{l,-l}}{P^{N}_{l-1,-(l-1)}}
    = \sqrt{\frac{2l(2l-1)}{(l-N)(l+N)}}\;\; s\,c,
  \label{eq:boundaryrecursion}
\end{equation}
valid for $l > |N|$, seeded at $l=|N|$ where the square root is one and
the value is $s^{2|N|}$ for $N\geq 0$ and $c^{2|N|}$ for $N<0$.  Both
boundaries ride on the same recursion, because the ratio depends on $N$
only through $N^2$ and because $P^N_{ll} = (-1)^{N+l}P^{-N}_{l,-l}$.

\subsection{Why the recursion is the safer form}

The closed form \eqref{eq:minorder} is evaluated in logarithms, and it has
to be: $\sqrt{(2l)!/\cdots}$ is of order $4^l$ and overflows long before
the value it belongs to, which is bounded by one.  The recursion
\eqref{eq:boundaryrecursion} multiplies by a factor of order $2sc =
\sin\theta \leq 1$ at each step, so nothing large is ever formed and no
$\log\Gamma$ is needed.  It is more robust than what it replaced, not
less.

Measured against the closed forms over $l_{\max}=128$, $|N|\leq 3$ and six
colatitudes, the worst relative discrepancy is $2\times10^{-13}$ in double
and $9\times10^{-17}$ in long double --- very nearly the same multiple of
$\varepsilon$, about 1000, in both, and growing linearly in $l_{\max}$
rather than with the number of entries compared.  That says the drift
belongs to the \emph{closed form}, which exponentiates a logarithm of size
$O(l\log 4)$ and loses bits in proportion.  The recursion's values are the
better ones, and the test's tolerance is written as $20\,l_{\max}
\varepsilon$ for that reason.

\subsection{The poles, and a data race}
\label{sec:signgam}

At $\theta=0$ and $\theta=\pi$ the ratio \eqref{eq:boundaryrecursion} is
$0\times\infty$.  No special case is needed: the seed is already the exact
$0$ or $1$ that the closed form's boundary branches return, and the factor
is zero thereafter, so the recursion produces the right answer unaided.

The seed row \eqref{eq:seedrow} was moved off the closed form for a
different reason.  glibc's \code{lgamma} writes the global
\code{signgam}, and the table is filled from every thread, so the closed
form is a data race there --- benign, since nothing reads \code{signgam},
but the only genuine one the library contains.  Generating the binomial
row instead removes the last $\log\Gamma$ from any threaded path, and the
values are exact rather than merely accurate.

Two things follow that are worth recording.  Grid construction is about
8\% faster, which is what the arithmetic predicts: the boundary is two
values per degree out of $2l+1$, so however expensive a $\log\Gamma$ is,
it is being paid on under 1\% of the table.  And ThreadSanitizer cannot
confirm the absence of the race, because GCC's \code{libgomp} carries no
TSan annotations and every OpenMP region is reported as racing with
whatever ran before it; a Clang build against a \code{libomp} compiled
with \code{LIBOMP\_TSAN\_SUPPORT} is the only way to get a real answer.

\section{Storage and access}
\label{sec:storage}

\subsection{The table}

Values are held as a flat array indexed by $(N, i_\theta)$ then by
$(l,m)$, contiguous in $(l,m)$ for each $(N,i_\theta)$ --- so a degree's
row is a contiguous run, which is what the transform's inner loop wants.
Offsets are closed form (\code{GSHIndices::OffsetForDegree}), not looked
up.

The table is large.  At $l_{\max}=256$ with $n_{\max}=2$ and all upper
indices it is $5\times 257\times(257^2-N^2)$ doubles, about 648\,MB; at
$l_{\max}=512$ it is 5.4\,GB.  This is why the grid is a value-semantic
handle over shared immutable state: copying one is a pointer copy, and
$\code{sizeof(Grid)}$ is 16 bytes.  Before that change, twenty copies of a
$l_{\max}=128$ grid cost 1638\,MB; afterwards, nothing.

\subsection{The row supplier}

Both transform loops take the row for a degree from an accessor rather
than walking a single iterator across the whole block.  The only thing
either loop needs is \emph{a contiguous run of values in $(l,m)$ order for
one $(N,i_\theta)$}, and asking for it one degree at a time costs a few
integer operations against a loop of length $(2l+1)k$.

That seam exists so the values can come from somewhere other than a table.
A second implementation runs the recursion into per-thread scratch inside
the transform, and a grid can be asked for it at construction.  It builds
no table at all: construction at $l_{\max}=256$ goes from 0.248\,s and
648\,MB to 0.015\,s and nothing.

\paragraph{It is slower, and the measurement is worth keeping.}  On an
8-core laptop the generated path is 1.4--3.2$\times$ slower than the
stored one in every configuration measured.  Batching narrows the gap,
since generation amortises over a chunk exactly as a table stream does,
but \emph{threads do not}: the ratio at eight threads is the ratio at one.
The premise that the stored path is pinned at the DRAM roof while a
generated one scales with cores is not visible on eight cores, and the
deployment target has more memory bandwidth per core, not less.  What
survives is the memory saving, which is large, and a NUMA argument that
only a two-socket machine can settle.

\section{Batching, chunking and threading}
\label{sec:batching}

\subsection{The batch}

The transform's primitive is ``transform $k$ fields sharing a grid, a
degree and an upper index'', with the single field as $k=1$.  What is
being amortised is the Wigner block for that upper index, which is why
fields at different upper indices cannot batch together --- for a rank-2
tensor that means batching over radii within each $N$, not across
components.

A batch is described by $(\text{count}, \text{stride}, \text{dist})$
rather than by contiguity: element $j$ of field $k$ lives at
$j\cdot\text{stride} + k\cdot\text{dist}$.  This is FFTW's advanced-interface
form, and it costs nothing at the stage that actually reads the caller's
layout, so it covers both arrangements the field layer produces:
\begin{center}
\begin{tabular}{lrr}
  \toprule
  & stride & dist \\
  \midrule
  components stored component by component & 1 & FieldSize \\
  components stored point by point & $n_{\text{components}}$ & 1 \\
  \bottomrule
\end{tabular}
\end{center}

\paragraph{The descriptor's uniformity is a constraint on tensor storage},
which was not obvious in advance.  Its members must be evenly spaced, and
in flat multi-index order the stored components sharing one upper index
are scattered at no fixed spacing as soon as there is any symmetry.  So a
tensor's buffer is ordered \emph{by upper index}, which makes each group a
contiguous run on both the field side and the coefficient side.  Ordered
any other way, the batching the primitive exists for would be unreachable
from the layer it was built for.

\paragraph{What batching is worth.}  About $2.3\times$ at the optimum,
not $k$.  The gain saturates because the stage is bound by per-element
load/store work rather than by table bandwidth, and batching removes
$(k-1)/k$ of the \emph{table} reads while leaving that untouched.  It also
\emph{reverses} beyond an optimum: the coefficient array is $k$ times
larger, and once it stops fitting in cache it is streamed from DRAM too.

\subsection{Chunking}

Because of that reversal, a batched call processes a large batch in chunks
of an internally chosen size rather than handing the caller's $k$ straight
to the inner loop.  The chunk is
\begin{equation}
  \text{chunk} \approx \frac{\text{cache}}
    {2 \times \text{copies} \times \text{bytes per coefficient block}},
  \label{eq:chunk}
\end{equation}
rounded to nearest, the factor of two being headroom for the streamed
Wigner row and the FFT output.

\paragraph{``Copies'' is not the thread count in both directions}, and
getting that wrong cost a factor of two.  The forward transform gives every
thread a private accumulator of $\text{chunk}\times\text{coefficientSize}$,
so its copies are its threads; the inverse gathers one block, shared and
read-only, so it has one copy however many threads read it.  Serving both
with the thread count starved the inverse --- at $l_{\max}=256$, $k=8$ on
eight threads it returned a chunk of one where the whole batch fits, and
the correction measured $2.0\times$ on that direction.

One thing the rule still cannot see: on a multi-CCD or multi-socket
machine a shared block is pulled into each cache domain that touches it,
so the inverse's single copy is really one per domain.

\subsection{Threading}

Threading is an explicit per-call policy and defaults to sequential.  The
library creates threads because it was asked to and not because it can,
and the rule it keeps is that \textbf{exactly one level threads}: an
operation asked to run in parallel from inside an existing parallel region
runs sequentially instead, so a caller parallelising over slices,
components or realisations cannot nest with the transform, and neither can
the table build.

The decomposition is over colatitudes, and the two directions differ.  The
inverse transform's colatitudes write disjoint rows and share only
read-only input, so they divide with no reduction at all.  The forward
transform's colatitudes all contribute to every coefficient, so each
thread accumulates into a private buffer and the partial sums are added at
the end --- partitioned rather than serialised: each thread sums every
thread's partials for its own block of coefficients.

\paragraph{Measured}, at $n=2$, complex, per call:
\begin{center}
\begin{tabular}{lrrrrr}
  \toprule
  & 1 thread & 2 & 4 & 8 & 16 \\
  \midrule
  $l_{\max}=128$, forward & 1.75\,ms & 1.89$\times$ & 2.98$\times$ & \textbf{4.02$\times$} & 3.13$\times$ \\
  $l_{\max}=128$, inverse & 1.77\,ms & 1.83$\times$ & 3.07$\times$ & 4.15$\times$ & \textbf{4.57$\times$} \\
  $l_{\max}=256$, forward & 12.16\,ms & 1.63$\times$ & 2.51$\times$ & \textbf{2.68$\times$} & 2.23$\times$ \\
  $l_{\max}=256$, inverse & 12.54\,ms & 1.69$\times$ & 2.64$\times$ & \textbf{2.98$\times$} & 2.97$\times$ \\
  \bottomrule
\end{tabular}
\end{center}

Sixteen threads is not better than eight and is often worse: the machine
has eight cores and sixteen hardware threads, and for memory-bound work
the second thread on a core adds contention rather than throughput.
Callers should ask for cores.

At $l_{\max}=128$ on eight threads the stage reaches 39--40\,GB/s against
a measured machine roof near 42, so it is genuinely bandwidth-bound there
and nothing further can be won without reducing the traffic itself.

\section{What the type system enforces}
\label{sec:types}

The field algebra is built so that a statement which is not mathematically
meaningful does not compile.  What follows is the list, with the reason in
each case.

\subsection{The index algebra}

Every node --- an owning field, a view, or an expression --- carries its
upper index as a compile-time constant, and the operators enforce
\S\ref{sec:components}:

\begin{center}
\begin{tabular}{lll}
  \toprule
  expression & upper index & requires \\
  \midrule
  $f+g$, $f-g$ & $N$ & equal upper indices \\
  $fg$ & $N_f + N_g$ & --- \\
  $f/g$ & $N_f$ & $N_g = 0$ \\
  $\conj{f}$ & $-N_f$ & --- \\
  $|f|$, $|f|^2$ & $0$ & --- \\
  $\mathrm{Re}\,f$, $\mathrm{Im}\,f$, $F(f)$ & $0$ & $N_f = 0$ \\
  $\int f \dOmega$ & --- & $N_f = 0$ \\
  \bottomrule
\end{tabular}
\end{center}

Three of these were wrong in the layer this replaced, and they are worth
naming.  \textbf{Conjugation reverses the upper index}: a quantity
carrying $e^{-iN\psi}$ has a conjugate carrying $e^{+iN\psi}$.
\textbf{Real and imaginary parts exist only at $N=0$}, because for
$N\neq 0$ the split is not preserved by a frame rotation and neither part
is a component of anything.  \textbf{Nonlinear functions likewise}, and
integration with them, since $\int f\dOmega$ vanishes identically unless
$N=0$.

The constraints are \code{requires}-clauses on the free operators, so an
unlawful combination is an overload-resolution failure at the call site
rather than an error inside a node --- which also means the negative cases
can be tested.

\subsection{Real-valuedness}

A node may be real-valued \emph{only} at $N=0$.  Real-valuedness is not
preserved by the frame rotation, so it is not a property any component of
any tensor can have elsewhere; a library able to represent one could
represent something that does not exist.  The constraint is closed under
every operation in the algebra --- $\mathrm{Re}$, $\mathrm{Im}$, $|f|$,
$|f|^2$ and $F(f)$ all land at $N=0$ precisely because that is the only
place their results could be real --- so it is checked once, on the
concept, and never re-derived.

\subsection{Laziness, and the aliasing theorem}

Operators build expression nodes; nothing is evaluated until a value is
read.  Every node in the layer is \textbf{pointwise and index-preserving}:
the value at $(i_\theta,i_\phi)$ reads its operands only at
$(i_\theta,i_\phi)$.  Hence
\begin{quote}
  $u = \text{expr}$ is safe when $u$ occurs in $\text{expr}$, and needs no
  temporary.
\end{quote}
That is a theorem about the node set rather than a property of any one
operator, which is why the invariant is stated in the stronger
``index-preserving'' form: a re-indexing view --- a longitude shift, a
transpose of the grid --- would be pointwise in the loose sense and would
break it.  Operations that are genuinely not pointwise (gradients,
raising and lowering, radial derivatives) live in the spectral layer for
this reason and not by accident.

Operand storage follows the value category: a terminal passed as an lvalue
is held by reference, everything else by value.  The residual hazard --- an
lvalue terminal destroyed while an expression referring to it is alive ---
is the same one Eigen has and is not removable in C++.

\subsection{Which components a tensor stores}

Both reductions of \S\ref{sec:reality} are computed by one \code{constexpr}
walk of the group generated by the permutation symmetry together with, for
a real tensor, the negation.  Reaching a component twice by different
routes is the interesting case rather than a failure: equating the two
routes produces a constraint.
\begin{center}
\begin{tabular}{ll}
  \toprule
  the two routes differ by & the orbit is \\
  \midrule
  a sign, same conjugation & identically zero \\
  conjugation, signs equal & real \\
  conjugation, signs opposite & purely imaginary \\
  \bottomrule
\end{tabular}
\end{center}
The first arises without reality at all --- it is the vanishing diagonal
of an antisymmetric tensor.  The others are the self-paired components,
which always sit at $N=0$, and are therefore exactly where a real-valued
field is admissible.

One consequence reaches the interface: a component accessor cannot return
one type.  A stored component is a view, a symmetric-permutation relative
is the same view, an antisymmetric relative is that view negated, a
reality relative is its conjugate at the reversed upper index, and a
vanishing orbit has nothing to return.  Traversal over all $3^p$
components is therefore a compile-time loop, not a runtime one.

\subsection{The two domains are not symmetric}

On the spatial side a derived component is pointwise --- $T^{-\alpha} =
(-1)^N\conj{T^{\alpha}}$ --- so a view plus an expression node gives it.
On the spectral side the corresponding relation is \eqref{eq:complevel},
$T^{-N}_{l,-m} = (-1)^m \conj{T^{N}_{lm}}$, which \emph{reverses the order
index}.  That is not a view over anything, so a tensor expansion exposes
only its stored components and deriving is done in the spatial domain.
This is the one place where the mirror between field and expansion breaks.

\section{Validation}
\label{sec:validation}

What each oracle actually pins, since a test suite that all passes is
uninformative about which facts are load-bearing.

\begin{center}
\begin{tabular}{lp{9.2cm}}
  \toprule
  oracle & what it pins \\
  \midrule
  \code{CheckWignerConvention}
    & That the stored value is $X^N_{lm}$ with the upper index first,
      against all nine $l=1$ values of D\&T (C.115).  Discriminating:
      the four entries with $N-m$ odd distinguish $d^l_{Nm}$ from
      $d^l_{mN}$. \\
  \code{CheckLegendre}
    & The $N=0$ row against \code{std::sph\_legendre}, which fixes the
      orthonormal scaling and the Condon--Shortley phase. \\
  \code{CheckAdditionTheorem}
    & \eqref{eq:addition} over all $|N|,|N'|\leq 40$ at $l_{\max}=40$, to
      $1000\varepsilon$.  The strongest single check on the recursion:
      boundary values feed the two-term recursion at every higher degree,
      so an error anywhere propagates into it. \\
  \code{CheckWignerBoundary}
    & The boundary recursion and the binomial seed row against the closed
      forms they replaced, to $20 l_{\max}\varepsilon$ --- a tolerance
      argued from where the drift actually is
      (\S\ref{sec:recursion}). \\
  \code{CheckCoeff2Coeff}
    & Coefficient-to-coefficient round trips. \\
  \code{TestRealFieldSymmetry}
    & \eqref{eq:basiclevel} at $n=2$ through complex transforms, and the
      rejection of real transforms at nonzero upper index. \\
  batched transform tests
    & A batch of $k$ against $k$ separate calls, at \emph{exact}
      equality.  Batching widens the inner loop without reordering any
      sum, so any difference at all would mean it had been restructured
      rather than widened. \\
  generated-path tests
    & The generating supplier against the stored table, again at exact
      equality: same recursion, same order, same rounding, only the
      storage differs. \\
  reality tests
    & \eqref{eq:reality} at every component of every tensor, and that the
      stored set costs $3^p$ reals per point. \\
  $\eth$ tests
    & \eqref{eq:ethidentities}, both of them. \\
  \bottomrule
\end{tabular}
\end{center}

\paragraph{On measuring this library.}  The development machine's clock
scales under load, and the noise floor between back-to-back runs of
identical code is about 10\%.  Comparisons must interleave the two
versions in one process; figures from different sessions are worth less
still.  Within-run ratios --- the stage split, the batching curve --- are
unaffected, because their terms are measured moments apart.

\section{The GEMM restructure}
\label{sec:gemm}

This was the one substantial piece of performance work the library had not
done, and \code{core-plan.md} \S10 named it as the lever aimed at what
actually binds.  \textbf{It is now built} --- \code{TransformKernel::Matrix()},
\code{core-plan.md} \S11, steps M1 to M6.

What follows is the argument as it was made, in enough detail to be argued
with, because the argument is what the measurement was read against; \S12.7
records what happened, including the two places where the prediction was
wrong.  The conditional voice below is deliberate and is not an oversight.

\subsection{The Legendre stage is a matrix product}

Written as it is implemented, the forward Legendre stage loops over
colatitudes and accumulates:
\begin{equation}
  f^{n}_{lm} \mathrel{+}= w_i\, X^{n}_{lm}(\theta_i)\, F_m(\theta_i)
  \qquad\text{for each } i, l, m .
  \label{eq:currentloop}
\end{equation}
Reordered so that the sum over $\theta$ is innermost, and at fixed upper
index $n$, it is one matrix product \emph{per order}:
\begin{equation}
  f^{n}_{lm} = \sum_{i} D^{(n,m)}_{l i}\, b^{(m)}_{i},
  \qquad
  D^{(n,m)}_{l i} = X^{n}_{lm}(\theta_i),
  \qquad
  b^{(m)}_{i} = w_i F_m(\theta_i).
  \label{eq:gemv}
\end{equation}
Over a batch of $k$ fields the right-hand side becomes a matrix and this is
a GEMM,
\begin{equation}
  \underbrace{(n_L(m) \times n_\theta)}_{D^{(n,m)}}
  \times
  \underbrace{(n_\theta \times k)}_{\text{weighted FFT output}}
  \;\longrightarrow\;
  \underbrace{(n_L(m) \times k)}_{\text{coefficients}},
  \label{eq:gemmshape}
\end{equation}
with $n_L(m) = l_{\max} - \max(|n|,|m|) + 1$ degrees surviving at that
order.  The inverse transform is the same matrix applied the other way
round, $(n_\theta \times n_L) \times (n_L \times k)$, so \emph{one stored
matrix serves both directions} through a transpose flag.

At $l_{\max}=256$ there are $2l_{\max}+1 = 513$ such products per upper
index, of heights $n_L$ from 257 down to 1, all with $n_\theta = 257$
columns.

\subsection{Why it should be faster}

Not for the reason P2 originally gave.  The stage is not short of
arithmetic intensity against DRAM; it is short of \emph{arithmetic per
memory operation} (\S\ref{sec:transform}).  Count them.

In \eqref{eq:currentloop}, one stored value $X$ is real and one FFT bin and
one accumulator are complex, so each visit costs, in doubles, one load of
$X$, two of the bin, and a load and a store of the accumulator --- seven
accesses --- for two fused multiply-adds, or four flops.  That is about
$0.6$ flops per double touched, and no amount of rewriting the loop moves
it, which is what the three hand-optimised variants of P1 found.

A GEMM microkernel holds a tile of the output in registers and streams the
two operands past it.  With a tile of $m_r$ rows by $n_r$ columns, each
step of the $K$ loop loads $m_r + n_r$ doubles and performs $m_r n_r$
fused multiply-adds:
\begin{equation}
  \frac{\text{flops}}{\text{double touched}}
    = \frac{2 m_r n_r}{m_r + n_r}
    \;\approx\; 5.3 \quad\text{at } m_r=4,\; n_r=8,
  \label{eq:microkernel}
\end{equation}
and the accumulator never leaves the registers, so its traffic vanishes
entirely rather than being counted.  That is an order of magnitude better
than the loop it replaces, on precisely the metric that binds.

\paragraph{The DRAM traffic is unchanged.}  The table is still read once
per batch; nothing about the restructure reduces it.  So the GEMM buys
nothing at all in the regime where the stage is bandwidth-bound --- which
it is, unbatched, on eight threads at $l_{\max}=128$, where it already
reaches 39--40\,GB/s against a roof near 42.  It buys everything in the
batched regime, where the same stage runs at about a fifth of the roof and
under a tenth of compute peak.  Since phases 2 to 5 always batch, that is
the regime that matters, but the distinction should be kept in view: a
benchmark run unbatched will show the restructure doing nothing, and be
right.

\subsection{What has to change}

\paragraph{The Wigner layout, which is the whole of step G's first option.}
\eqref{eq:gemmshape} needs $D^{(n,m)}$ contiguous in $(l,\theta)$ for fixed
$(n,m)$.  The table is currently contiguous in $(l,m)$ for fixed
$(n,\theta)$, which is the transpose of what is wanted.  Storing it as
$[n][m][l][\theta]$ makes each matrix a contiguous row-major block with
leading dimension $n_\theta$, and the total size is unchanged --- summing
$n_L(m)\, n_\theta$ over $m$ gives the same $66\,045 \times 257$ doubles
per upper index.  \code{Wigner}'s existing \code{Storage} tag orders the
$(n,\theta)$ axes only; this is a finer and different axis order, and it is
what \code{RowMajor}'s deferred fate ([C7]) was waiting for.

\paragraph{The FFT output has to land in $m$-major order}, since
\eqref{eq:gemmshape} wants, for one $m$, a contiguous $(\theta,\kappa)$
block.  That means all the FFTs run \emph{before} any of the Legendre
stage, rather than interleaved per colatitude as now --- the intermediate
is $n_\phi \times n_\theta \times k$ complex, which is just the field, 2.1\,MB
at $l_{\max}=256$ and $k=1$.

The transpose is nearly free, for the same reason it was free in tier 1
(P6): a \code{plan\_many} with output stride equal to the number of
transforms and output distance 1 writes the result directly in
$[m][\theta][\kappa]$ order.  \emph{Nearly}, and this paragraph used to say
\emph{free} and to specify a single call over all $n_\theta k$ rows.  Both
were corrected by measurement at \code{core-plan.md} M2.  The strided write
costs four to eight times as much when the number of transforms times 16 is a
power of two, because successive orders then collide in the same cache sets;
and a \emph{small} block of colatitudes per call beats one large one by
1.7--2$\times$ while wanting eighty times less workspace, since a strided
write over a few hundred bytes stays inside a couple of cache lines and one
over tens of kilobytes does not.

\paragraph{The GEMM is real, and that is a gift.}  $D$ is real while the
data are complex, so there is no standard complex-times-real BLAS call ---
but there does not need to be.  A complex $(n_\theta \times k)$ block with
$\kappa$ fastest \emph{is} a real $(n_\theta \times 2k)$ block, since
\code{std::complex} stores its parts adjacently.  So the whole stage is
\code{dgemm} with $N = 2k$, and the doubling helps most exactly where $k$
is small.

\subsection{What it does to threading}

The products at different $m$ write disjoint outputs.  So the natural
decomposition is over orders, and it needs \textbf{no accumulator and no
reduction in either direction} --- which is why \S10 says this restructure
subsumes [C11].  The forward transform's thread-private accumulators, which
are correct at eight threads and are predicted not to survive 64--128,
simply stop existing.

Two cautions.  The work per order is not constant: $n_L(m)$ falls linearly
in $|m|$, so a static split over orders is roughly twice as unbalanced as
it looks, and the schedule should divide work rather than count.  And a
thread's share of the table is a set of whole $(n,m)$ blocks, which is a
better NUMA story than the current colatitude split --- first touch during
the table build can be made to match the split exactly, which it currently
cannot.

\subsection{What to expect, honestly}

The shape is unfavourable.  A GEMM reaches a large fraction of peak when
all three dimensions are large; here $M = n_L$ is large only for small
$|m|$, $K = n_\theta$ is a few hundred, and $N = 2k$ is between 2 and 16.
Skinny-$N$ GEMMs are the case general BLAS kernels handle worst, and the
513 products per upper index are individually small enough that call
overhead is not negligible.  A batched BLAS interface, where one call takes
the whole set, is the right shape and not universally available.

So the honest expectation is a factor of two or three in the batched
regime, not the order of magnitude that \eqref{eq:microkernel} suggests in
isolation --- and the measurement is the point of doing it.

\subsection{What it composes with, and what it fights}

\paragraph{The reflection composes, and helps more here than it would
now.}  D\&T (C.118) gives $P^{N}_{lm}(-\mu) = (-1)^{l+N} P^{N}_{l,-m}(\mu)$,
so the matrices at $+m$ and $-m$ are related by reversing the colatitude
and a sign that alternates with $l$.  Splitting the input into its
even and odd parts about $\theta = \pi/2$ turns the pair of products at
$\pm m$ into two products of half the height, halving both the stored
table and the arithmetic.  This is the standard trick in production
spherical-harmonic libraries.  It is also exactly the reduction that P4
warned would trade bandwidth for an irregular access pattern in the
\emph{current} loop; inside a GEMM the same reduction is a clean halving of
$K$, so the objection does not carry over.

\paragraph{It fights on-the-fly generation, and not merely by ordering.}
The recursion runs up degrees at fixed $(n,\theta)$, producing values in
exactly the order the \emph{present} layout stores them and exactly the
wrong order for $[n][m][l][\theta]$.  It would be easy to read that as an
inconvenience with a transpose as its price.  It is not.  The recursion's
output for one $(n,\theta)$ spans \emph{every order at once}, so isolating
the single $(n,m)$ block that a per-order product wants means either keeping
all of it --- which is the table, and not having one is the entire purpose of
generation --- or re-running the recursion once per order, $2l_{\max}+1$
times over.  There is no third way, so the two are not combinable rather than
merely awkward to combine, and \code{core-plan.md} [C17] refuses the pair at
grid construction.  Since the generated path measured slower anyway
(\S\ref{sec:storage}), nothing is lost.

\paragraph{It costs the exact-equality tests, and buys a better one.}  A
GEMM sums in whatever order its kernel chooses, so a GEMM path cannot be
bit-compared against the current one; the batched-against-unbatched tests,
which presently demand exact equality precisely because tier 1 does
\emph{not} reorder any sum, become tolerance comparisons on that path.
\code{core-plan.md} [C12] turns that trade around by keeping \emph{both}
kernels permanently: the loop path is then an oracle for the matrix path on
identical inputs, which checks the layout, the indexing, the FFT ordering
and the accumulation together, and which does not exist today.  The two are
not independent in the $d$ values, since both read the same recursion ---
but that half is pinned separately, against the $l=1$ table and against
\code{std::sph\_legendre}.

\paragraph{Which BLAS.}  An earlier draft of this paragraph said Eigen was
already a dependency and could supply the kernel.  It is not: Eigen left the
tree when GaussQuad dropped it, and nothing in the library names it now.  So
there is no free GEMM to start from, and \code{core-plan.md} [C14] takes the
decision rather than defaulting into one --- an optional BLAS, found and not
fetched, under \code{GSHTRANS\_WITH\_BLAS}, in the pattern the optional
\code{Interpolation} dependency established.  A build without it does not
offer the matrix kernel, exactly as a build without \code{Interpolation}
does not offer \code{SplineDerivative}.  It would be this library's first
non-header-only dependency, which is why it is an option and not a
requirement.

\subsection{What it measured, and where the argument was wrong}
\label{sec:gemmdone}

Built as \code{core-plan.md} \S11 describes, in six steps: the transposed
table, the Fourier stage run first and landed order-major, the forward and
inverse kernels, threading over orders, the measurement, and the reflection.
Batched at $l_{\max}=256$ on eight threads it is worth \textbf{5.8 to
6.0$\times$ forward and 2.9 to 3.2$\times$ inverse}, and it stores half the
table.

\paragraph{The honest estimate was right, and low.}  \S\ref{sec:gemm} put it
at a factor of two or three batched.  Sequentially the forward gives
$3.0\times$, inside that; threaded it gives $4.6\times$, outside it.  The
excess is not the GEMM doing better but the loop kernel doing worse: at that
size the loop forward scales only $1.5\times$ from one thread to eight,
because eight private accumulators of 8.4\,MB compete for a 16\,MB cache.
The matrix kernel has no accumulator to collapse.  So the threaded comparison
measures the GEMM \emph{and} the deletion of that problem together.

\paragraph{The prediction that nothing is available unbatched was confirmed,
and refined.}  At $l_{\max}=256$, $k=1$ on eight threads the answer is
$1.20\times$ forward and $1.10\times$ inverse, which is nothing, and it is the
only regime in which that is true.  The refinement is the reason: this
document said no gain was available because \emph{the loop kernel} is already
at the roof.  Measured, the loop kernel is at 65\% of it and the matrix kernel
at 100\%, which is why a small gain survives.  The conclusion stands and is
now demonstrated rather than predicted.

\paragraph{The reflection halves the table and not the arithmetic.}
\S\ref{sec:gemm} says it halves both.  Only the table halves: the same
$2l_{\max}+1$ products are still done, of the same shapes.  Converting the
symmetry into arithmetic --- pairing $\pm m$ into one product of double the
width, which is the natural reading --- was built and \emph{measured worse}
sequentially, so it was taken out and only the halving kept.  What the
reflection buys is traffic and memory, which is what this document's own DRAM
argument should have predicted.

\paragraph{And the free transpose is free only away from one hazard.}  The
output write has stride \code{howMany} complex doubles, so when
$16\,\code{howMany}$ is a power of two the writes for successive orders
collide in the same cache sets, and the stage costs four to eight times what
it costs at \code{howMany} $\pm 2$.  The Fourier stage is therefore blocked
over a few colatitudes rather than run in one call, which is faster anyway
and wants eighty times less workspace.

\section{Interpolation}
\label{sec:interpolation}

\code{Interpolate(field, scheme)} returns a callable of $(\theta,\phi)$
modelling \code{ScalarFunctionS2}, which is what a field constructor takes,
so remeshing is one line.  Three schemes: \code{Scheme::Spectral()} sums the
expansion and is exact for a band-limited field at $O(l_{\max}^2)$ a point;
\code{Scheme::Bilinear()} and \code{Scheme::Bicubic()} are
\code{Interpolation}'s rectilinear schemes over the samples, $O(1)$ a point
and second and fourth order respectively.

Two facts about the sphere make this more than a forwarding call.  The
longitudes run $0,\dots,2\pi-\Delta\phi$, so the last cell has nothing to
interpolate against; and neither pole is a grid point, so a local scheme
extrapolates there --- exactly where $d^l_{Nm}\sim\sin^{|m|}\theta$ is most
delicate.  Both are fixed by handing the scheme a \emph{padded} grid: a wrap
column, which is exact since $\phi=2\pi$ is $\phi=0$, and two polar rows.

The polar rows are computed rather than guessed.  At a pole all but one order
vanishes:
\begin{equation}
  f(0,\phi) = e^{+iN\phi}\sum_{l} f^{N}_{l,+N}\sqrt{\tfrac{2l+1}{4\pi}},
  \qquad
  f(\pi,\phi) = e^{-iN\phi}\sum_{l} (-1)^{l-N} f^{N}_{l,-N}
                 \sqrt{\tfrac{2l+1}{4\pi}} .
  \label{eq:polerows}
\end{equation}
The $\phi$ dependence is the point: a spin-weighted field at a coordinate
pole is not single-valued, because $\ebasis{\pm}$ depends on the azimuth of
approach, so a polar row is a \emph{row} and a constant one would be wrong at
every $N\neq 0$.  Computing it needs the expansion, which is why building a
local interpolant costs a forward transform and why its cheapness is per
evaluation rather than per interpolant.

Measured, the padding works: the polar cells and the last longitude cell come
out two to five times \emph{better} than the interior at every resolution, so
the regions the padding was built for are not the weak ones.  A local scheme
wants an oversampled grid --- bicubic carries 13\% error at the band limit,
$7\times10^{-4}$ at four times it and $4\times10^{-5}$ at eight --- and
\code{ForBand} is how that room is asked for.  One whole-grid remesh costs
about what fifty spectral point evaluations cost, so past a few dozen
scattered points the composite of the two beats either alone.

\section{The 3-j symbols}
\label{sec:threej}

\code{3j.h} tabulates the whole $(m_1,m_3)$ plane at fixed degrees, with
$m_2=-(m_1+m_3)$ fixed by the selection rule.  The table is the primitive: a
single symbol costs one.

The values come from the Schulten--Gordon recursion.  The order recurrence
\begin{equation}
  A(m_2)\,g(m_2) + B(m_2)\,g(m_2-1) + A(m_2-1)\,g(m_2-2) = 0
  \label{eq:threejrec}
\end{equation}
has a growing and a decaying solution, with a classically allowed region
between two forbidden ones.  Recursing \emph{inward} from a forbidden end
follows the growing solution, so rounding error decays; outward follows the
decaying one, so it grows exponentially.  Each row is therefore built inward
from both ends, matched where the halves overlap, and normalised from the
unitary property $(2l_1+1)\sum_{m_2} g(m_2)^2 = 1$.  No closed-form seed is
needed and no factorial is ever formed.

Two details make it practical.  The turning point is found by watching the
recurrence coefficient --- while $|c_1|$ decreases the run is heading towards
larger values and is stable, and the first increase says to come from the
other end --- rather than by locating it analytically.  And the phase is
\emph{tracked} through the match rather than read back from
$g(m_{2\max})$: a near-stretched row at high degree spans a dynamic range of
$10^{201}$, so the rescaling flushes that element to zero and the sign with
it.

This replaced two earlier schemes, and \code{docs/3j-plan.md} records the
measurements.  A one-directional recursion fails exponentially near
\emph{stretched} triangles --- where one degree approaches the sum of the
other two, which is the top of every coupling sum rather than an exotic
corner --- and Racah's closed form covers exactly that region and no more.
Together they still left a band uncovered; Schulten--Gordon covers all of it,
at 1.1 to 1.3 times the cost.

\paragraph{What checks it.}  Not the completeness relation, which the
algorithm normalises by and which therefore holds by construction; and not the
recurrence residual alone, which is homogeneous and so blind to an overall
scale or sign.  The check that sees everything is invariance under a cyclic
permutation of the columns, which runs the recursion along entirely different
lines: it agrees to $10^{-16}$ at $(200,200,200)$ and $10^{-15}$ at
$(1000,1000,1999)$.  Racah is kept in the test tree as an independent formula,
restricted to entries where its alternating sum is short and it is exact.
\eqref{eq:threejrec} is checked at run time, so a table that has gone wrong is
refused rather than returned.

\section{Choosing a policy by measurement}
\label{sec:tuning}

Several of the library's choices are properties of the machine rather than of
the mathematics, and \code{core-plan.md} \S12 records that they cannot be
settled by reasoning.  \code{Tuning.h} times the alternatives on the caller's
own problem and returns \emph{values} --- never a configured grid, so nothing
is substituted silently, which matters most for the piece of machinery most
likely to do so.

The chunking policy and the planner flag were moved off the shared
implementation onto the handle, since neither decides the table; \code{With()}
changes either for a pointer copy, and the result shares \code{Identity()}
with its parent, so fields are interchangeable between them.  Without that a
sweep of four candidate chunks rebuilt the table four times to choose an
integer.

Measured on one laptop: the kernel choice is worth 1.9 to $5.8\times$ and
picks the matrix kernel in every configuration tried, while the chunk is worth
at most $1.28\times$ and nothing at all in thirteen of eighteen.  Both are
cheap enough to run at start-up, which is why no persistence layer was built:
the measurement is cheaper than remembering it, and \code{core-plan.md}
\S12.6 names the size at which that reverses.  A candidate must beat the
incumbent by ten per cent to displace it, that being the measured noise floor.

\section{What has not been done, and why}
\label{sec:whatnext}

Recorded so that the absences read as decisions.

\paragraph{The symmetry reductions of \S\ref{sec:legendre}} would cut the
table fourfold.  They are not implemented because in the regime that
matters they would buy little: batched, with a chunk that fits, the
Legendre stage runs at about a fifth of the memory roof and under a tenth
of compute peak, so halving the traffic addresses neither binding
constraint.  Unbatched they would be worth roughly $2\times$; batching
already took that.

\paragraph{Reduced-precision storage} is in the same position, and for the
same reason.  It also needs a tighter round-trip oracle than exists before
the experiment would mean anything.

\paragraph{Polar truncation}, skipping the $(m,\theta)$ pairs where
$d^l_{Nm}$ is negligible --- it falls off like $\sin^{|m|}\theta$ --- is the
one option that makes the problem smaller rather than making the machine work
harder.  \emph{This paragraph used to add that it probably dominates both of
the reductions above.  Measured, it does not.}  \code{core-plan.md} \S11.7
prices it at about $1.15\times$ at $l_{\max}=256$ for the rectangular trim a
GEMM can take, and $1.40\times$ at $l_{\max}=1024$ for the triangular ceiling
--- against the $1.5$ to $2\times$ that was assumed.  The decay is real but
slower to a useful tolerance than the argument suggests: at $10^{-15}$,
$\sin^{256}\theta$ is still above threshold over a third of the range.  It is
dropped rather than deferred.

\paragraph{The wisdom store} of \S\ref{sec:tuning} --- a keyed,
fingerprinted file remembering what was measured --- is not built.  Both
tuners are cheap enough to run at start-up, so there is nothing expensive
enough to be worth remembering; FFTW's wisdom is worth persisting because
planning can take minutes, and this takes about a second.

\paragraph{Generating the $d$-values inside the transform} exists
(\S\ref{sec:storage}) and its deeper forms do not.  Fusing the recursion with
the consumer, and blocking it over colatitudes, would halve its arithmetic;
they are not built because the first rung already measured slower than the
stored path on every configuration of the development machine, and halving
the arithmetic cannot close that gap.  What the path is for is memory, and
that it delivers.

\paragraph{The 6-j symbols} are not implemented.  The path is known and is
the same one: Schulten and Gordon give the 6-j recursion in the same paper as
the 3-j one, and \S\ref{sec:threej}'s machinery --- inward from both ends,
match, normalise from a unitary property --- transfers directly.

\paragraph{The conventions of \S\ref{sec:remains}} were open when this
document was first written and are now closed, jointly, by reading a gradient
back into the physical frame.  What is left there is one convention rather
than two unknowns, and a test observes it.

\end{document}
